\documentclass[final,3p,times]{elsarticle}

\usepackage[T1]{fontenc}
\usepackage[utf8]{inputenc}
\usepackage{lmodern}
\usepackage{amsmath,amssymb,amsthm,mathtools}
\usepackage{graphicx}
\usepackage{booktabs}
\usepackage{tikz}
\usepackage{hyperref}

\graphicspath{{figures/}{./}}

%% theorem environments (common counter)
\theoremstyle{plain}
\newtheorem{theorem}{Theorem}
\newtheorem{lemma}[theorem]{Lemma}
\newtheorem{proposition}[theorem]{Proposition}
\theoremstyle{definition}
\newtheorem{definition}[theorem]{Definition}
\theoremstyle{remark}
\newtheorem{remark}[theorem]{Remark}

%% compact notation for the method name
\newcommand{\dziti}{\ensuremath{\delta}\text{-Ziti}}
\newcommand{\supp}{\operatorname{supp}}
\newcommand{\spanop}{\operatorname{span}}
\newcommand{\dd}{\mathrm{d}}

\begin{document}

\begin{frontmatter}

\title{Theoretical Analysis and Numerical Investigations of the \dziti{} Method
for Three-Dimensional Transient Diffusion Problems}

\author{Rajae Malek}
\affiliation{organization={Business School},
	addressline={Mundiapolis University, Honoris United Universities},
	city={Casablanca},
	country={Morocco}}

\begin{abstract}
We analyse the \dziti{} method for three-dimensional transient diffusion with
Neumann boundary conditions. The method is built from compactly supported
mollifiers, orthonormalised by a two-term Gram--Schmidt recurrence that is
justified by the tridiagonal structure of the associated Gram matrix. The
resulting basis is used in two complementary ways: as an $L^2$ projection
space, and through an interpolatory reconstruction at the roots of the basis
functions, which yields an explicit fully discrete scheme without numerical
quadrature or inversion of a mass matrix. We prove that $f_h$ is the best least-squares approximation in the
\dziti{} space and that $\lVert f-f_h\rVert_{L^2}\le\lVert f-I_h f\rVert_{L^2}$
with $I_h f=\sum_i f(x_i)\varphi_i(x)$. A second-order bound holds for the
mollification $\eta_h*f$, not for $f_h$: the space $V_h$ does not contain
constants, so no estimate $\lVert f-f_h\rVert_{L^2}\le C h^2\lVert f''\rVert_{L^\infty}$
can hold for all $f\in W^{2,\infty}$. The semi-discrete Galerkin scheme is unconditionally
energy-stable. The corresponding explicit Euler scheme is
stable under a CFL restriction of order $\Delta t=\mathcal{O}(h^2)$. The
construction is extended to three dimensions by tensorisation on Cartesian
domains and by a spherical-coordinate (or chord-sweeping) strategy on the unit
ball. Numerical tests on the unit cube and the unit ball, including
manufactured solutions, are used to assess accuracy and computational cost
against standard linear finite elements.
\end{abstract}

\begin{keyword}
\dziti{} method \sep orthonormal basis \sep Gram--Schmidt process \sep
explicit scheme \sep collocation \sep transient diffusion \sep heat equation
\sep convergence analysis \sep energy stability
\end{keyword}

\end{frontmatter}

%==============================================================================
\section{Introduction}
\label{sec:intro}

The heat equation is the prototype of linear parabolic diffusion. On a bounded
Lipschitz domain $\Omega\subset\mathbb{R}^d$ ($d=3$ in the applications below)
and a time interval $(0,T]$, it reads
\begin{equation}
\label{eq:heat}
\begin{cases}
\dfrac{\partial u}{\partial t}-\nabla\cdot\bigl(k\nabla u\bigr)=f
& \text{in }\Omega\times(0,T],\\[4pt]
k\dfrac{\partial u}{\partial n}=g
& \text{on }\partial\Omega\times(0,T],\\[4pt]
u(\cdot,0)=u_0
& \text{in }\Omega,
\end{cases}
\end{equation}
where $u$ is the unknown temperature, $k\ge k_{\min}>0$ is the thermal
conductivity (taken constant in the analysis), $f$ is an internal heat source,
$g$ is a prescribed boundary flux, $n$ is the outward unit normal, and $u_0$ is
the initial temperature. The Neumann condition models insulated or
flux-specified surfaces. Under standard assumptions
$f\in L^2(0,T;L^2(\Omega))$, $g\in L^2(0,T;H^{-1/2}(\partial\Omega))$ and
$u_0\in L^2(\Omega)$, problem~\eqref{eq:heat} admits a unique weak solution
$u\in L^2(0,T;H^1(\Omega))$ with $\partial_t u\in L^2(0,T;(H^1(\Omega))')$
\cite{Evans2010,Brezis2011}. Accurate discretisations remain demanding in three
dimensions in the presence of steep gradients or boundary layers, because
stability of explicit time stepping forces $\Delta t=\mathcal{O}(h^2)$ while
low-order spatial schemes require fine meshes to control the error.

Classical discretisations include finite differences on simple geometries
\cite{Gibou2007} and finite elements on unstructured meshes
\cite{Babuska1973,Quarteroni2008,Thomee1997}. Linear finite elements are
robust, but they typically require substantial refinement near regions of high
curvature of $u$, which is costly for transient three-dimensional computations
\cite{Mohamed2013,OHara2009}. Enrichment strategies such as the partition of
unity finite element method (PUFEM) \cite{Melenk1996}, the generalised finite
element method (GFEM) \cite{Strouboulis2000} and the discontinuous enrichment
method \cite{Farhat2001,Mohamed2014} alleviate this by inserting
problem-adapted functions into the trial space. In particular,
Malek~\emph{et al.}~\cite{Malek2019} developed a PUFEM scheme for
three-dimensional transient diffusion, using time-independent exponential
enrichments and thereby reducing the number of degrees of freedom relative to
standard finite elements.

The \dziti{} method, introduced by Bsiss and Ziti
\cite{Bsiss2016,Bsiss2017,Bsiss2023} and subsequently extended by Malek and
Ziti \cite{Malek2021,Malek2019Stokes,Malek2022}, follows a different route. A
family of compactly supported mollifiers is orthonormalised in $L^2$, and the
roots of the resulting basis are used as collocation nodes. The scheme
therefore starts from a Galerkin identity and, after replacing integrals by
point values at those roots, behaves like a meshless/particle method. It has
been applied to hyperbolic conservation laws, including Riemann problems for
Burgers' equation \cite{Bsiss2016,Yoon2008}, and to elliptic and parabolic
problems on Cartesian domains and on the disk \cite{Malek2021}. A distinctive
feature, inherited from the generating kernel, is that the basis functions
approximate a Dirac mass in the sense of distributions, which is useful for
detecting loss of regularity \cite{Malek2022,Malek2023}.

In the present work we apply the method to~\eqref{eq:heat} in three space
dimensions and put the spatial approximation and the time-stepping scheme on a
rigorous footing. The contributions are as follows.
\begin{itemize}
\item We recall the one-dimensional construction and give a complete
justification of the two-term Gram--Schmidt recurrence from the compact
overlap of the generating functions.
\item We prove that the reconstruction $f_h$ is the best least-squares
approximation in $V_h$ and that $\lVert f-f_h\rVert_{L^2}\le\lVert f-I_h f\rVert_{L^2}$
with $I_h f=\sum_{i}f(x_i)\varphi_i(x)$. The kernel $\eta_h$ is a second-order
approximate identity, but $V_h$ does not contain constants, and therefore
no bound $\lVert f-f_h\rVert_{L^2}\le C h^2\lVert f''\rVert_{L^\infty}$ holds
for all $f\in W^{2,\infty}$.
We also establish energy stability of the semi-discrete Galerkin scheme and a
CFL condition for the explicit Euler method.
\item We extend the scheme to the unit cube by tensorisation and to the unit
ball by spherical coordinates (and, alternatively, by sweeping chords), and we
write the explicit three-dimensional Neumann scheme.
\end{itemize}

The rest of the paper is organised as follows. Section~\ref{sec:method}
presents the construction of the \dziti{} basis, the interpolatory formulae
and the three-dimensional extension. Section~\ref{sec:theory} contains the
theoretical results. Section~\ref{sec:numerics} describes the numerical
experiments. Conclusions are given in Section~\ref{sec:concl}.

%==============================================================================
\section{The \dziti{} method}
\label{sec:method}

\subsection{Generating kernel}
\label{sec:kernel}

The construction starts from the standard $C^\infty_c$ bump
\begin{equation}
\label{eq:Phi}
\Phi(x)=
\begin{cases}
\exp\bigl(1/(\lvert x\rvert^2-1)\bigr) & \text{if }\lvert x\rvert<1,\\
0 & \text{otherwise.}
\end{cases}
\end{equation}
Equivalently, $\Phi(x)=\exp\bigl(-1/(1-\lvert x\rvert^2)\bigr)$ on the open
unit interval. The function $\Phi$ is even, nonnegative and compactly
supported in $(-1,1)$. Set
\begin{equation}
\label{eq:C}
C:=\Biggl(\int_{\mathbb{R}}\Phi(x)\,\dd x\Biggr)^{-1}
\simeq 2.252283621,
\end{equation}
and, for $\varepsilon>0$,
\begin{equation}
\label{eq:mollifier}
\eta_\varepsilon(x):=\frac{C}{\varepsilon}\,\Phi\Bigl(\frac{x}{\varepsilon}\Bigr).
\end{equation}
Then $\int_{\mathbb{R}}\eta_\varepsilon=1$, $\eta_\varepsilon$ is even, and
$\eta_\varepsilon\to\delta_0$ in $\mathcal{D}'(\mathbb{R})$ as
$\varepsilon\to 0$. This Dirac approximation is the origin of the prefix
``$\delta$'' in the name of the method \cite{Bsiss2016,Malek2021}.

\subsection{One-dimensional generators and orthonormalisation}
\label{sec:1d}

Let $I=[a,b]\subset\mathbb{R}$ and let $m\ge 2$ be an integer. Set
$h:=(b-a)/m$ and $x_i:=a+(i-1)h$ for $i=1,\dots,m+1$, so that $x_1=a$ and
$x_{m+1}=b$. The generators are the translates
\begin{equation}
\label{eq:phi}
\varphi_i(x):=\eta_h(x-x_i)=\frac{C}{h}\,\Phi\Bigl(\frac{x-x_i}{h}\Bigr),
\qquad i=1,\dots,m+1,
\end{equation}
restricted to $I$. Consequently
\begin{equation}
\label{eq:supports}
\begin{aligned}
\supp(\varphi_1)&=[x_1,x_2],\\
\supp(\varphi_i)&=[x_{i-1},x_{i+1}] && \text{for }i=2,\dots,m,\\
\supp(\varphi_{m+1})&=[x_m,x_{m+1}].
\end{aligned}
\end{equation}
In particular, $\varphi_i$ and $\varphi_j$ have supports of disjoint interior
whenever $\lvert i-j\rvert\ge 2$, so that
$\langle\varphi_i,\varphi_j\rangle_{L^2(I)}=0$ for $\lvert i-j\rvert\ge 2$.
The family $(\varphi_i)_{1\le i\le m+1}$ is linearly independent (the Gram
matrix is symmetric positive definite and tridiagonal). Representative
profiles are shown in Figure~\ref{fig:phi}.

\begin{figure}[htbp]
\centering
\begin{tikzpicture}
\node[inner sep=0pt] {\includegraphics[width=0.85\textwidth]{Fig2}};
\node at (-4.5,3) {\tiny $\varphi_{1}$};
\node at (-1,3) {\tiny $\varphi_{i-1}$};
\node at (0.1,3) {\tiny $\varphi_{i}$};
\node at (1.2,3) {\tiny $\varphi_{i+1}$};
\node at (4.5,3) {\tiny $\varphi_{m+1}$};
\end{tikzpicture}
\caption{Generators $(\varphi_i)$ on $I=[-4,4]$ with $m=8$
($C\simeq 2.252283621$).}
\label{fig:phi}
\end{figure}

Let $\langle\cdot,\cdot\rangle$ denote the inner product of $L^2(I)$. Because
the Gram matrix of $(\varphi_i)$ is tridiagonal, the Gram--Schmidt process
reduces to a two-term recurrence: the orthogonal family
$(\widetilde\psi_i)_{1\le i\le m+1}$ is defined by
\begin{equation}
\label{eq:GS}
\begin{cases}
\widetilde\psi_1:=\varphi_1,\\[4pt]
\widetilde\psi_i:=\varphi_i+\lambda_{i-1}\widetilde\psi_{i-1},
& i=2,\dots,m+1,\\[4pt]
\lambda_{i-1}
:=-\dfrac{\langle\varphi_i,\widetilde\psi_{i-1}\rangle}
{\langle\widetilde\psi_{i-1},\widetilde\psi_{i-1}\rangle}.
\end{cases}
\end{equation}
Unfolding the recurrence yields
\begin{equation}
\label{eq:unfold}
\widetilde\psi_i
=\varphi_i+\lambda_{i-1}\varphi_{i-1}
+\lambda_{i-1}\lambda_{i-2}\varphi_{i-2}
+\cdots
+\Biggl(\prod_{k=1}^{i-1}\lambda_k\Biggr)\varphi_1,
\end{equation}
and therefore
$\supp(\widetilde\psi_i)=[x_1,x_{i+1}]$ for $i\le m$ (with the obvious
modification for $i=m+1$). The orthonormal basis is then
\begin{equation}
\label{eq:psi}
\psi_i:=\frac{\widetilde\psi_i}{\lVert\widetilde\psi_i\rVert_{L^2(I)}},
\qquad i=1,\dots,m+1.
\end{equation}
Consecutive profiles are shown in Figure~\ref{fig:psi}. We write
$V_h:=\spanop\{\psi_i:1\le i\le m+1\}=\spanop\{\varphi_i:1\le i\le m+1\}$.

\begin{figure}[htbp]
\centering
\begin{tikzpicture}
\node[inner sep=0pt] {\includegraphics[width=0.85\textwidth]{Fig3}};
\node at (-1,2.1) {\tiny $\psi_{i-1}$};
\node at (0.1,2.1) {\tiny $\psi_{i}$};
\node at (1.2,2.1) {\tiny $\psi_{i+1}$};
\end{tikzpicture}
\caption{Consecutive orthonormal functions $\psi_{i-1}$, $\psi_i$ and
$\psi_{i+1}$ on $I=[-4,4]$ with $m=8$.}
\label{fig:psi}
\end{figure}

The coefficients $(\lambda_i)$ obey the nonlinear recurrence
\begin{equation}
\label{eq:lambda-rec}
\lambda_1=-\frac{\langle\varphi_1,\varphi_2\rangle}{\langle\varphi_1,\varphi_1\rangle},
\qquad
\lambda_{i+1}=\frac{\lambda_1}{2-\lambda_1\lambda_i},
\qquad i\ge 1,
\end{equation}
which follows from the two-term structure and the translation invariance of
the interior generators \cite{Bsiss2023,Malek2022}. One has
$-1<\lambda_i<0$. Moreover, $(\lambda_i)$ converges geometrically to a limit
$\lambda_\star\in(-1,0)$. Quantitatively, for every $\varepsilon>0$,
\begin{equation}
\label{eq:N0}
\lvert\lambda_{i+1}-\lambda_i\rvert<\varepsilon
\quad\text{as soon as}\quad
i\ge N_0
:=\Biggl\lfloor
\frac{\displaystyle\ln\biggl(\frac{\varepsilon(2-\lambda_1^2)}{\lambda_1^3-\lambda_1}\biggr)}
{\displaystyle\ln\biggl(\Bigl(\frac{\lambda_1}{2+\lambda_1}\Bigr)^2\biggr)}
\Biggr\rfloor+1.
\end{equation}
In practice $N_0=9$ already yields stationarity at machine precision for
$\varepsilon=10^{-9}$; see Table~\ref{tab:lambda}. From that rank onward the
recurrence need not be recomputed, which is the source of the method's
algorithmic compression \cite{Bsiss2023}.

\begin{table}[htbp]
\centering
\caption{Values of $\lambda_i$ produced by~\eqref{eq:GS}--\eqref{eq:lambda-rec}.
The sequence is stationary at $\lambda_\star\simeq -0.273518325$ from the rank
$N_0=9$.}
\label{tab:lambda}
\begin{tabular}{lccccc}
\toprule
$i$ & $1$ & $2$ & $3$ & $4$ & $5$ \\
\midrule
$\lambda_i$
& $-0.508960181$ & $-0.292344636$ & $-0.274934054$
& $-0.273624280$ & $-0.273526252$ \\
\midrule
$i$ & $6$ & $7$ & $8$ & $N_0=9$ & $10$--$101$ \\
\midrule
$\lambda_i$
& $-0.273518918$ & $-0.273518369$ & $-0.273518329$
& $-0.273518325$ & $-0.273518325$ \\
\bottomrule
\end{tabular}
\end{table}

\subsection{Roots and interpolatory reconstruction}
\label{sec:roots}

The following localisation of the roots is proved in
\cite{Malek2022,Malek2021,Bsiss2017}. We restate it in a form that also
defines the collocation nodes used below.

\begin{proposition}[Roots of the orthonormal family]
\label{prop:roots}
Let $m\ge 2$. There exists a unique increasing sequence
$(r_k)_{1\le k\le m+1}$ in $I$ such that
\begin{enumerate}
\item $r_k\in(x_k,x_{k+1})$ for $k=1,\dots,m$, and $r_{m+1}=b$;
\item for every $i\in\{1,\dots,m+1\}$,
$\psi_i(r_i)>0$ and $\psi_i(r_k)=0$ whenever $k\neq i$.
\end{enumerate}
In particular, for $i\ge 2$, the function $\psi_i$ has exactly $i-1$ roots in
$(x_1,x_i)$, namely $r_1,\dots,r_{i-1}$. The remaining identities
$\psi_i(r_k)=0$ for $k>i$ follow from the support
$\supp(\widetilde\psi_i)\subset[x_1,x_{i+1}]$.
\end{proposition}

The biorthogonality $\psi_i(r_k)=\delta_{ik}\psi_i(r_i)$ is the key algebraic
property of the method. Given $f\in C^0(I)$, one therefore has a uniquely
defined interpolant in $V_h$,
\begin{equation}
\label{eq:interp}
I^\delta f:=\sum_{i=1}^{m+1}c_i[f]\,\psi_i,
\qquad
c_i[f]:=\frac{f(r_i)}{\psi_i(r_i)},
\end{equation}
which satisfies $I^\delta f(r_k)=f(r_k)$ for all $k$. Equivalently, if a
function is expanded as $f\simeq\sum_i c_i\psi_i$ and the expansion is
evaluated at $x=r_k$, the off-diagonal terms vanish and one recovers
$c_k=f(r_k)/\psi_k(r_k)$. This is an interpolatory identity; it is \emph{not}
the formula for the $L^2$ Fourier coefficients
$\langle f,\psi_i\rangle$, except when $f\in V_h$.

Integrating the interpolant gives a consistent quadrature,
\begin{equation}
\label{eq:quad}
\int_a^b I^\delta f(x)\,\dd x
=\sum_{i=1}^{m+1}\frac{f(r_i)}{\psi_i(r_i)}\,I_i,
\qquad
I_i:=\int_a^b\psi_i(x)\,\dd x.
\end{equation}
If one further approximates $I_i$ by the one-point rule
$I_i\simeq 1/\psi_i(r_i)$, one obtains the simplified formula
\begin{equation}
\label{eq:quad-simple}
\int_a^b f(x)\,\dd x
\simeq\sum_{i=1}^{m+1}\frac{f(r_i)}{\psi_i(r_i)^2},
\end{equation}
which appears in earlier presentations of the method. Throughout the analysis
we use~\eqref{eq:interp} and~\eqref{eq:quad}, which do not rely on that extra
approximation. Likewise,
\begin{equation}
\label{eq:inner-interp}
\langle I^\delta f,\psi_i\rangle
=\frac{f(r_i)}{\psi_i(r_i)},
\end{equation}
so the identity $\langle f,\psi_i\rangle\simeq f(r_i)/\psi_i(r_i)$ is exact on
$V_h$ and is used as a quadrature for general $f$.

\begin{proposition}[Interpolatory formulae]
\label{prop:interp}
Let $m\ge 2$ and $f\in C^0(I)$. The operator $I^\delta$ defined
by~\eqref{eq:interp} is the unique interpolant in $V_h$ at the nodes
$(r_i)$, and~\eqref{eq:quad}--\eqref{eq:inner-interp} hold. In particular,
\begin{equation}
\label{eq:ci}
c_i[f]=\frac{f(r_i)}{\psi_i(r_i)},
\qquad
\int_a^b f(x)\psi_i(x)\,\dd x
\simeq\frac{f(r_i)}{\psi_i(r_i)}.
\end{equation}
\end{proposition}

\subsection{Tensorisation in two and three dimensions}
\label{sec:tensor}

On a rectangle $\Omega=[a,b]\times[c,d]$ one constructs one-dimensional
families $(\psi_i^{[1]})_{1\le i\le m_1+1}$ on $[a,b]$ and
$(\psi_j^{[2]})_{1\le j\le m_2+1}$ on $[c,d]$, with collocation nodes
$(r_i)$ and $(s_j)$ respectively. The tensor-product basis
\begin{equation}
\label{eq:2d-basis}
\psi_{ij}(x,y):=\psi_i^{[1]}(x)\,\psi_j^{[2]}(y)
\end{equation}
is orthonormal in $L^2(\Omega)$. The interpolant of $f\in C^0(\Omega)$ is
\begin{equation}
\label{eq:2d-interp}
I^\delta f(x,y)
=\sum_{i=1}^{m_1+1}\sum_{j=1}^{m_2+1}
\frac{f(r_i,s_j)}{\psi_i^{[1]}(r_i)\,\psi_j^{[2]}(s_j)}
\,\psi_{ij}(x,y),
\end{equation}
and the quadrature analogue of~\eqref{eq:quad} reads
\begin{equation}
\label{eq:2d-quad}
\int_a^b\int_c^d I^\delta f
=\sum_{i=1}^{m_1+1}\sum_{j=1}^{m_2+1}
\frac{f(r_i,s_j)}{\psi_i^{[1]}(r_i)\,\psi_j^{[2]}(s_j)}
\,I_i^{[1]}\,I_j^{[2]}.
\end{equation}

On the unit cube $\Omega=[a,b]^3$ the same construction produces an
orthonormal tensor basis
\begin{equation}
\label{eq:3d-basis}
\psi_{ijk}(x,y,z)
:=\psi_i(x)\,\psi_j(y)\,\psi_k(z),
\qquad
1\le i,j,k\le m+1,
\end{equation}
with collocation nodes $(r_i,r_j,r_k)$ and
\begin{equation}
\label{eq:3d-interp}
I^\delta f
=\sum_{i,j,k=1}^{m+1}
\frac{f(r_i,r_j,r_k)}{\psi_i(r_i)\psi_j(r_j)\psi_k(r_k)}
\,\psi_{ijk}.
\end{equation}
We write $V_h(\Omega):=\spanop\{\psi_{ijk}\}$ and denote by
$\mathcal{R}_h:=\{(r_i,r_j,r_k)\}$ the three-dimensional collocation grid.
The number of degrees of freedom is $N_h=(m+1)^3$.

On the unit ball $B:=\{x\in\mathbb{R}^3:\lvert x\rvert<1\}$ two strategies
are available, following \cite{Malek2021}. The first sweeps $B$ by chords
parallel to the coordinate axes and repeats the one-dimensional construction
on each chord, with a fixed number of nodes per chord and a chord-dependent
step. The second uses spherical coordinates
$x=\rho\sin\theta\cos\phi$, $y=\rho\sin\theta\sin\phi$,
$z=\rho\cos\theta$, with
$(\rho,\theta,\phi)\in[0,1]\times[0,\pi]\times[0,2\pi]$, and the tensor
product of three one-dimensional \dziti{} families on those intervals. The
volume element $\rho^2\sin\theta\,\dd\rho\,\dd\theta\,\dd\phi$ is
incorporated into the quadrature~\eqref{eq:quad}. Both strategies reduce
every integral and every collocation identity to the one-dimensional
formulae of Section~\ref{sec:roots}.

%==============================================================================
\section{Theoretical analysis}
\label{sec:theory}

\subsection{Algebraic properties of the basis}

\begin{lemma}[Equality of spans and best approximation]
\label{lem:orthonormal}
Let $I=[a,b]\subset\mathbb{R}$, and let the two families
$\{\varphi_i\}_{i=1}^{m+1}$ and $\{\psi_i\}_{i=1}^{m+1}$ be those constructed
by the \dziti{} method in Section~\ref{sec:1d}. Set
$V_h:=\spanop\{\psi_i\}_{i=1}^{m+1}$. Then:
\begin{enumerate}
\item The two families generate the same vector space:
\begin{equation}
\label{eq:spans}
\spanop\{\psi_i\}_{i=1}^{m+1}
=
\spanop\{\varphi_i\}_{i=1}^{m+1}.
\end{equation}
\item In the continuous case, and likewise in the discrete case, the
\dziti{} reconstruction
\begin{equation}
\label{eq:fh}
f_h(x)
=\sum_{i=1}^{m+1}\frac{f(r_i)}{\psi_i(r_i)}\,\psi_i(x)
\end{equation}
is the best approximation of $f$ in $V_h$ in the least-squares sense. Consequently
\begin{equation}
\label{eq:best}
\lVert f-f_h\rVert_{L^2(I)}
=\inf_{v\in V_h}\lVert f-v\rVert_{L^2(I)}
\le
\lVert f-I_h f\rVert_{L^2(I)},
\end{equation}
where $I_h f$ is the interpolant constructed from the generators
$\{\varphi_i\}_{i=1}^{m+1}$,
\begin{equation}
\label{interp}
I_h f=\sum_{i=1}^{m+1}f(x_i)\,\varphi_i(x).
\end{equation}
\end{enumerate}
In~\eqref{eq:best} the norm is that of $L^2(I)$ in the continuous case, and
the discrete least-squares norm associated with the roots $(r_i)$ in the
discrete case.
\end{lemma}

\begin{proof}
(i) The inclusion
$\spanop\{\psi_i\}\subset\spanop\{\varphi_i\}$ follows from the definition of
the Gram--Schmidt family: each $\widetilde\psi_i$ (and hence each $\psi_i$)
is a linear combination of $\varphi_1,\dots,\varphi_i$. The reverse inclusion
is obtained by a simple recurrence from~\eqref{eq:GS},
\begin{equation}
\label{eq:inverse}
\varphi_1=\widetilde\psi_1,
\qquad
\varphi_i=\widetilde\psi_i-\lambda_{i-1}\widetilde\psi_{i-1},
\quad i=2,\dots,m+1.
\end{equation}
Thus each $\varphi_i$ is a linear combination of $\widetilde\psi_i$ and
$\widetilde\psi_{i-1}$. Since $\psi_i$ is a nonzero multiple of
$\widetilde\psi_i$, one has~\eqref{eq:spans}.

(ii) \textbf{Continuous case.}
We equip $L^2(I)$ with the inner product
\begin{equation}
\label{eq:ip}
(f,g)=\int_I f(x)g(x)\,\dd x.
\end{equation}
The family $\{\psi_i\}_{i=1}^{m+1}$ is an orthonormal basis of $V_h$, so every
$v\in V_h$ may be written
\begin{equation}
v(x)=\sum_{i=1}^{m+1}a_i\psi_i(x).
\end{equation}
Consider the quadratic functional
\begin{align*}
E(a_1,\dots,a_{m+1})
&=\lVert f-v\rVert_{L^2(I)}^2 \\
&=(f,f)-2\sum_{i=1}^{m+1}a_i(\psi_i,f)+\sum_{i=1}^{m+1}a_i^2 \\
&=(f,f)-2\sum_{i=1}^{m+1}a_i c_i+\sum_{i=1}^{m+1}a_i^2,
\end{align*}
where $c_i$ is the \dziti{} coefficient of $f$. By the interpolatory formulae
of Proposition~\ref{prop:interp},
\begin{equation}
\label{eq:ci-id}
c_i
=(\psi_i,f)
=\int_a^b f(x)\psi_i(x)\,\dd x
=\frac{f(r_i)}{\psi_i(r_i)}.
\end{equation}
A critical point of $E$ is characterised by the vanishing of all partial
derivatives:
\begin{equation}
\frac{\partial E}{\partial a_i}=-2c_i+2a_i.
\end{equation}
The condition $\partial E/\partial a_i=0$ therefore yields $a_i=c_i$. The
Hessian of $E$ is twice the identity, so this critical point is the unique
global minimiser. Hence
\begin{equation}
f_h(x)=\sum_{i=1}^{m+1}\frac{f(r_i)}{\psi_i(r_i)}\,\psi_i(x)
\end{equation}
is the best approximation of $f$ in $V_h$ in the $L^2(I)$ least-squares sense,
and
\begin{equation}
\lVert f-f_h\rVert_{L^2(I)}
=\inf_{v\in V_h}\lVert f-v\rVert_{L^2(I)}.
\end{equation}
From (i), the interpolant $I_h f$ defined by~\eqref{interp} belongs to $V_h$.
Taking $v=I_h f$ in the infimum gives the comparison with~\eqref{interp}:
\begin{equation}
\lVert f-f_h\rVert_{L^2(I)}
\le
\lVert f-I_h f\rVert_{L^2(I)}.
\end{equation}

\textbf{Discrete case.}
Let the inner product be the discrete sum associated with the \dziti{} roots,
\begin{equation}
\label{eq:discrete-ip}
(f,g)_h
:=\sum_{k=1}^{m+1}\frac{f(r_k)\,g(r_k)}{\psi_k(r_k)^2}.
\end{equation}
The biorthogonality $\psi_i(r_k)=\delta_{ik}\psi_i(r_i)$ of
Proposition~\ref{prop:roots} implies
$(\psi_i,\psi_j)_h=\delta_{ij}$ and
\begin{equation}
(f,\psi_i)_h=\frac{f(r_i)}{\psi_i(r_i)}.
\end{equation}
Repeating the argument above with $(\cdot,\cdot)_h$ in place of~\eqref{eq:ip},
the same functional $E$ is minimised if and only if
$a_i=(f,\psi_i)_h=f(r_i)/\psi_i(r_i)$. Thus $f_h$ is exactly the discrete
least-squares projection of $f$ onto $V_h$, and~\eqref{eq:best} holds for the
norm $\lVert\cdot\rVert_h$ induced by~\eqref{eq:discrete-ip}.
\end{proof}

\subsection{Error estimate}

Lemma~\ref{lem:orthonormal} gives the comparison
$\lVert f-f_h\rVert_{L^2(I)}\le\lVert f-I_h f\rVert_{L^2(I)}$ with
$I_h f=\sum_i f(x_i)\varphi_i$, but that comparison does not produce a rate.
The only second-order estimate attached to the construction is the
mollification bound for $\eta_h*f$. The reconstruction $f_h$ itself cannot
satisfy a bound of order $h^2$ in terms of $\lVert f''\rVert_{L^\infty}$,
because $V_h$ does not contain constants.

\begin{theorem}[Second-order mollification]
\label{thm:mollify}
Let $K(s):=C\Phi(s)$. Then $K$ is even, compactly supported in $(-1,1)$, and
$\int_{\mathbb{R}}K=1$. Set
$\mu_2:=\int_{\mathbb{R}}s^2 K(s)\,\dd s\in(0,\infty)$.
If $f\in W^{2,\infty}(\mathbb{R})$, then
\begin{equation}
\label{eq:mollify-inf}
\lVert\eta_h*f-f\rVert_{L^\infty(\mathbb{R})}
\le
\frac{\mu_2}{2}h^2\lVert f''\rVert_{L^\infty(\mathbb{R})},
\end{equation}
and therefore, on any bounded interval $I$,
\begin{equation}
\label{eq:mollify-L2}
\lVert\eta_h*f-f\rVert_{L^2(I)}
\le
\frac{\mu_2}{2}\lvert I\rvert^{1/2}h^2\lVert f''\rVert_{L^\infty(\mathbb{R})}.
\end{equation}
The same bounds hold for $f\in W^{2,\infty}(I)$ after a $W^{2,\infty}$
extension to an $h$-neighbourhood of $I$.
\end{theorem}

\begin{proof}
Evenness of $K$ gives $\int_{\mathbb{R}}s K(s)\,\dd s=0$. The Taylor formula
with integral remainder yields
\begin{equation}
f(x-hs)-f(x)
=-hs f'(x)
+\int_0^1(1-t)h^2 s^2 f''(x-ths)\,\dd t.
\end{equation}
Multiply by $K(s)$ and integrate in $s$. The linear term vanishes, and
\begin{equation}
\bigl\lvert(\eta_h*f)(x)-f(x)\bigr\rvert
\le
h^2\lVert f''\rVert_{L^\infty(\mathbb{R})}
\int_{\mathbb{R}}\int_0^1(1-t)s^2 K(s)\,\dd t\,\dd s
=\frac{\mu_2}{2}h^2\lVert f''\rVert_{L^\infty(\mathbb{R})},
\end{equation}
which is~\eqref{eq:mollify-inf}. Integrating the square of this bound over
$I$ gives~\eqref{eq:mollify-L2}.
\end{proof}

\begin{theorem}[No $O(h^2)$ bound for $f_h$ in $L^2$]
\label{thm:error}
Let $I=[a,b]$ and let $f_h$ be any function in $V_h$, in particular the
\dziti{} reconstruction~\eqref{eq:fh}. There is no constant $C$, independent
of $h$ and $f$, such that
\begin{equation}
\label{eq:err-false}
\lVert f-f_h\rVert_{L^2(I)}
\le
C\,\lvert I\rvert^{1/2}\,h^2\,\lVert f''\rVert_{L^\infty(I)}
\end{equation}
holds for every $f\in W^{2,\infty}(I)$ and every mesh size $h=(b-a)/m$.
Quantitatively, there exists $\gamma>0$, depending only on $\Phi$, such that
the constant function $f\equiv 1$ satisfies
\begin{equation}
\label{eq:err-lower}
\lVert 1-v\rVert_{L^2(I)}
\ge
\gamma\lvert I\rvert^{1/2}
\qquad\text{for every }v\in V_h\text{ and every }h.
\end{equation}
In particular~\eqref{eq:err-lower} holds for $v=f_h$.
\end{theorem}

\begin{proof}
On each cell $[x_k,x_{k+1}]$, $k=1,\dots,m$, at most $\varphi_k$ and
$\varphi_{k+1}$ are nonzero. For $x=x_k+\theta h$ with $\theta\in(0,1)$
and $v=\sum_i a_i\varphi_i\in V_h$,
\begin{equation}
v(x)
=\frac{C}{h}\bigl(a_k\Phi(\theta)+a_{k+1}\Phi(\theta-1)\bigr)
=\alpha_k\Phi(\theta)+\alpha_{k+1}\Phi(1-\theta),
\end{equation}
where $\alpha_i:=C a_i/h$ and evenness of $\Phi$ was used. Thus
\begin{equation}
\lVert 1-v\rVert_{L^2(I)}^2
=
\sum_{k=1}^{m}
h\int_0^1\bigl\lvert 1-\alpha_k\Phi(\theta)-\alpha_{k+1}\Phi(1-\theta)\bigr\rvert^2
\,\dd\theta.
\end{equation}
Set
\begin{equation}
\gamma^2
:=\inf_{\alpha,\beta\in\mathbb{R}}
\int_0^1\bigl\lvert 1-\alpha\Phi(\theta)-\beta\Phi(1-\theta)\bigr\rvert^2\,\dd\theta.
\end{equation}
The infimum is attained (the integrand is a positive definite quadratic in
$(\alpha,\beta)$ plus lower-order terms) and is strictly positive: the three
functions $1$, $\Phi(\theta)$ and $\Phi(1-\theta)$ are linearly independent in
$L^2(0,1)$. Indeed, if $A+B\Phi(\theta)+D\Phi(1-\theta)=0$ almost everywhere,
then evaluating at $\theta\to 0^+$, $\theta\to 1^-$ and $\theta=1/2$ gives
$A+B/e=0$, $A+D/e=0$ and $A+(B+D)\Phi(1/2)=0$, with
$\Phi(1/2)=\exp(-4/3)$. Hence $B=D$ and
$B\bigl(2\exp(-4/3)-e^{-1}\bigr)=0$. The factor in parentheses is nonzero
because $e^{1/3}\neq 2$. Thus $A=B=D=0$, so $\gamma>0$.

Therefore
$\lVert 1-v\rVert_{L^2(I)}^2\ge m h\gamma^2=\lvert I\rvert\gamma^2$,
which is~\eqref{eq:err-lower}.

If~\eqref{eq:err-false} held, the choice $f\equiv 1$ (so $f''=0$) would force
$f_h=1$ almost everywhere, contradicting~\eqref{eq:err-lower}.
\end{proof}

\begin{remark}
\label{rem:mollify-vs-fh}
Theorem~\ref{thm:mollify} is the precise sense in which the generating kernel
is second order: convolution against $\eta_h$ reproduces $W^{2,\infty}$ data
up to $O(h^2)$. That operator is not the reconstruction $f_h$, and
$\eta_h*f\notin V_h$ in general. Lemma~\ref{lem:orthonormal} still yields
$\lVert f-f_h\rVert_{L^2(I)}\le\lVert f-I_h f\rVert_{L^2(I)}$, but the
right-hand side does not tend to zero like $O(h^2)$, consistently with
Theorem~\ref{thm:error}.
\end{remark}

\subsection{Weak form and semi-discrete stability}

Multiplying~\eqref{eq:heat} by $v\in H^1(\Omega)$ and integrating by parts
gives the weak formulation: find
$u\in L^2(0,T;H^1(\Omega))$ such that, for a.e.\ $t\in(0,T)$,
\begin{equation}
\label{eq:weak}
\bigl(\partial_t u(t),v\bigr)
+a\bigl(u(t),v\bigr)
=\ell(t;v)
\qquad\forall v\in H^1(\Omega),
\end{equation}
where
\begin{equation}
\label{eq:forms}
a(w,v):=\int_\Omega k\nabla w\cdot\nabla v\,\dd x,
\qquad
\ell(t;v):=\int_\Omega f(t)v\,\dd x+\int_{\partial\Omega}g(t)v\,\dd\sigma.
\end{equation}
The bilinear form $a$ is continuous and coercive on $H^1(\Omega)$ up to the
kernel of the Neumann problem (constants); uniqueness of~\eqref{eq:heat} is
restored by the initial condition.

Let $\{\psi_\alpha\}_{\alpha\in\mathcal{I}_h}$ be the orthonormal \dziti{}
basis of $V_h(\Omega)\subset H^1(\Omega)$ constructed in
Section~\ref{sec:tensor} (tensor product on the cube, spherical or
chord-wise construction on the ball). The semi-discrete Galerkin scheme
consists in seeking $u_h(t)=\sum_{\alpha}c_\alpha(t)\psi_\alpha\in V_h(\Omega)$
such that $u_h(0)$ is the \dziti{} reconstruction~\eqref{eq:fh} of the
initial datum $u_0$, and
\begin{equation}
\label{eq:semidisc}
\bigl(\partial_t u_h(t),v_h\bigr)
+a\bigl(u_h(t),v_h\bigr)
=\ell(t;v_h)
\qquad\forall v_h\in V_h(\Omega).
\end{equation}
Orthonormality reduces the mass matrix to the identity: writing $C(t)=(c_\alpha(t))$,
\begin{equation}
\label{eq:ode}
C'(t)+AC(t)=F(t),
\end{equation}
where $A_{\alpha\beta}=a(\psi_\beta,\psi_\alpha)$ is the stiffness matrix and
$F_\alpha(t)=\ell(t;\psi_\alpha)$. No mass-matrix inversion is required.

\begin{theorem}[Energy stability of the semi-discrete scheme]
\label{thm:semi}
Let $k>0$ be constant and let $u_h\in C^1([0,T];V_h(\Omega))$
solve~\eqref{eq:semidisc}.
\begin{enumerate}
\item If $f=0$ and $g=0$, the scheme dissipates energy:
\begin{equation}
\label{eq:energy-hom}
\lVert u_h(t)\rVert_{L^2(\Omega)}^2
+2k\int_0^t\lVert\nabla u_h(s)\rVert_{L^2(\Omega)}^2\,\dd s
=
\lVert u_h(0)\rVert_{L^2(\Omega)}^2.
\end{equation}
\item If $f\in L^2(0,T;L^2(\Omega))$ and
$g\in L^2(0,T;L^2(\partial\Omega))$, there exists $C_T>0$, depending only
on $\Omega$, $k$ and $T$, such that for every $t\in[0,T]$,
\begin{equation}
\label{eq:energy}
\begin{aligned}
\lVert u_h(t)\rVert_{L^2(\Omega)}^2
&+k\int_0^t\lVert\nabla u_h(s)\rVert_{L^2(\Omega)}^2\,\dd s \\
&\le
C_T\Biggl(
\lVert u_h(0)\rVert_{L^2(\Omega)}^2
+\int_0^t\Bigl(
\lVert f(s)\rVert_{L^2(\Omega)}^2
+\lVert g(s)\rVert_{L^2(\partial\Omega)}^2
\Bigr)\,\dd s
\Biggr).
\end{aligned}
\end{equation}
\end{enumerate}
The estimates are independent of $h$. In particular, the semi-discrete
scheme is unconditionally stable in $L^2(\Omega)$.
\end{theorem}

\begin{proof}
Take $v_h=u_h(t)$ in~\eqref{eq:semidisc}. Since
$a(u_h,u_h)=k\lVert\nabla u_h\rVert_{L^2(\Omega)}^2$,
\begin{equation}
\label{eq:energy-id}
\frac12\frac{\dd}{\dd t}\lVert u_h\rVert_{L^2(\Omega)}^2
+k\lVert\nabla u_h\rVert_{L^2(\Omega)}^2
=(f,u_h)+\langle g,u_h\rangle_{\partial\Omega}.
\end{equation}
If $f=0$ and $g=0$, the right-hand side vanishes. Integrating in time
yields~\eqref{eq:energy-hom}.

In the inhomogeneous case, the Cauchy--Schwarz inequality and the trace
theorem give
\begin{equation}
\begin{aligned}
\bigl\lvert(f,u_h)+\langle g,u_h\rangle_{\partial\Omega}\bigr\rvert
&\le
\lVert f\rVert_{L^2(\Omega)}\lVert u_h\rVert_{L^2(\Omega)}\\
&\qquad
+C_{\mathrm{tr}}\lVert g\rVert_{L^2(\partial\Omega)}
\bigl(\lVert u_h\rVert_{L^2(\Omega)}+\lVert\nabla u_h\rVert_{L^2(\Omega)}\bigr).
\end{aligned}
\end{equation}
Young's inequality absorbs the gradient trace term into the dissipation:
\begin{equation}
\begin{aligned}
\lVert f\rVert\lVert u_h\rVert
&\le \tfrac12\lVert u_h\rVert_{L^2(\Omega)}^2+\tfrac12\lVert f\rVert_{L^2(\Omega)}^2,\\
C_{\mathrm{tr}}\lVert g\rVert\lVert u_h\rVert
&\le \tfrac12\lVert u_h\rVert_{L^2(\Omega)}^2+C_1\lVert g\rVert_{L^2(\partial\Omega)}^2,\\
C_{\mathrm{tr}}\lVert g\rVert\lVert\nabla u_h\rVert
&\le \tfrac{k}{2}\lVert\nabla u_h\rVert_{L^2(\Omega)}^2
+C_2\lVert g\rVert_{L^2(\partial\Omega)}^2,
\end{aligned}
\end{equation}
with $C_1,C_2$ depending only on $\Omega$ and $k$. Therefore
\begin{equation}
\label{eq:energy-diff}
\frac{\dd}{\dd t}\lVert u_h\rVert_{L^2(\Omega)}^2
+k\lVert\nabla u_h\rVert_{L^2(\Omega)}^2
\le
2\lVert u_h\rVert_{L^2(\Omega)}^2
+\lVert f\rVert_{L^2(\Omega)}^2
+C_3\lVert g\rVert_{L^2(\partial\Omega)}^2.
\end{equation}
Dropping the gradient term and applying Gronwall's lemma, we obtain
\begin{equation}
\lVert u_h(t)\rVert_{L^2(\Omega)}^2
\le
e^{2t}\lVert u_h(0)\rVert_{L^2(\Omega)}^2
+\int_0^t e^{2(t-s)}
\Bigl(\lVert f(s)\rVert_{L^2(\Omega)}^2
+C_3\lVert g(s)\rVert_{L^2(\partial\Omega)}^2\Bigr)\,\dd s.
\end{equation}
Integrating~\eqref{eq:energy-diff} over $[0,t]$ and inserting this bound on
$\int_0^t\lVert u_h\rVert_{L^2(\Omega)}^2\,\dd s$ produces~\eqref{eq:energy},
with $C_T$ depending on $\Omega$, $k$ and $T$ only.
\end{proof}

\begin{remark}[Forward Euler on the Galerkin ODE]
\label{rem:galerkin-euler}
Because the \dziti{} basis is orthonormal,~\eqref{eq:ode} has mass matrix
equal to the identity. Advancing it by forward Euler,
$C^{n+1}=(I-\Delta t\,A)C^n+\Delta t\,F^n$, is stable in $\ell^2$ if and
only if $\Delta t\,\rho(A)\le 2$. The stiffness matrix $A$ is symmetric
positive semi-definite, and the inverse inequality
$\lVert\nabla v_h\rVert_{L^2(\Omega)}\le C_{\mathrm{inv}}h^{-1}\lVert v_h\rVert_{L^2(\Omega)}$
on $V_h$ yields $\rho(A)\le C_{\mathrm{inv}}^2 k\,h^{-2}$, with
$C_{\mathrm{inv}}$ independent of $h$. A sufficient CFL condition is therefore
$k\Delta t\le 2h^2/C_{\mathrm{inv}}^2$. This restriction is of the same order
as~\eqref{eq:CFL} below; the constant $C_{\mathrm{inv}}$ is not needed for
the collocation scheme analysed next.
\end{remark}

\subsection{Fully discrete explicit scheme}

Let $\Delta t=T/N_t$ and $t_n=n\Delta t$. Write $u_\alpha^n$ for the
approximate value of $u$ at the collocation node $x_\alpha\in\mathcal{R}_h$
and at time $t_n$. Combining the interpolatory identities of
Proposition~\ref{prop:interp} with a second-order difference for the Laplacian
on the (in general nonuniform) nodes $(r_i)$ produces an explicit scheme. On
the unit cube, with $u_{ijk}^n\simeq u(t_n,r_i,r_j,r_k)$ and
$h_i:=r_{i+1}-r_i$, the standard second divided difference
\begin{equation}
\label{eq:dxx}
\delta_{xx}u_{ijk}^n
:=\frac{2}{h_i+h_{i-1}}
\Biggl(
\frac{u_{i+1,jk}^n-u_{ijk}^n}{h_i}
-\frac{u_{ijk}^n-u_{i-1,jk}^n}{h_{i-1}}
\Biggr)
\end{equation}
is used in each coordinate (and likewise $\delta_{yy}$, $\delta_{zz}$). The
interior update is forward Euler,
\begin{equation}
\label{eq:euler}
u_{ijk}^{n+1}
=u_{ijk}^n
+\Delta t\,k\,(\delta_{xx}+\delta_{yy}+\delta_{zz})u_{ijk}^n
+\Delta t\,f(t_n,r_i,r_j,r_k),
\end{equation}
for $2\le i,j,k\le m$. At the faces of the cube the Neumann condition is
enforced by a symmetric ghost value realising $k\partial_n u=g$; this keeps
the discrete Laplacian self-adjoint, which is used in the stability argument
below. A first-order one-sided closure may be used in computations, but it
need not be symmetric, and the $\ell^2$ contraction of Theorem~\ref{thm:scheme}
is then no longer guaranteed by the same energy identity. On the unit ball
the same time-stepping is applied to the spherical Laplacian
\begin{equation}
\label{eq:sph-lap}
\Delta u
=u_{\rho\rho}+\frac{2}{\rho}u_\rho
+\frac{1}{\rho^2\sin\theta}\partial_\theta(\sin\theta\,u_\theta)
+\frac{1}{\rho^2\sin^2\theta}u_{\phi\phi},
\end{equation}
with $\partial_\rho u=g/k$ at $\rho=1$ and the usual regularity conditions
at $\rho=0$ and on the polar axis.

\begin{theorem}[CFL stability of the explicit collocation scheme]
\label{thm:scheme}
Let $\Delta_h$ denote the standard seven-point discrete Laplacian on a
uniform Cartesian mesh of step $h$, with homogeneous Neumann (or periodic)
closure realised by symmetric ghost values. Then $\Delta_h$ is symmetric and
negative semi-definite on $\ell^2(\mathcal{R}_h)$, and its eigenvalues lie in
$[-\Lambda_h,0]$ with $\Lambda_h=12/h^2$. Write $\tau=k\Delta t$ and consider
the homogeneous forward Euler scheme
\begin{equation}
\label{eq:euler-hom}
U^{n+1}=U^n+\tau\Delta_h U^n=(I+\tau\Delta_h)U^n.
\end{equation}
The following are equivalent:
\begin{enumerate}
\item $k\Delta t\le h^2/6$;
\item $\lvert 1+\tau\lambda\rvert\le 1$ for every eigenvalue $\lambda$ of
$\Delta_h$;
\item $\lVert I+\tau\Delta_h\rVert_{\ell^2\to\ell^2}\le 1$.
\end{enumerate}
In particular, if
\begin{equation}
\label{eq:CFL}
k\Delta t\le\frac{h^2}{6},
\end{equation}
then
\begin{equation}
\label{eq:l2stab}
\lVert U^{n+1}\rVert_{\ell^2(\mathcal{R}_h)}
\le
\lVert U^n\rVert_{\ell^2(\mathcal{R}_h)},
\qquad
\max_{0\le n\le N_t}\lVert U^n\rVert_{\ell^2(\mathcal{R}_h)}
\le
\lVert U^0\rVert_{\ell^2(\mathcal{R}_h)}.
\end{equation}

On a nonuniform mesh, let $h_{\min}:=\min_i h_i$ and let
\begin{equation}
\label{eq:weighted}
(U,V)_h
:=\sum_{i,j,k}U_{ijk}V_{ijk}\,\omega_i\omega_j\omega_k,
\qquad
\omega_i:=\frac{h_i+h_{i-1}}{2},
\end{equation}
be the weighted inner product associated with the cell volumes. If the
stencil~\eqref{eq:dxx} is used with a symmetric Neumann closure, then
$\Delta_h$ is self-adjoint and negative semi-definite on this inner product.
Gershgorin's theorem yields $\Lambda_h\le 12/h_{\min}^2$, and the same energy
argument gives the sufficient condition
\begin{equation}
\label{eq:CFL-nu}
k\Delta t\le\frac{h_{\min}^2}{6}
\end{equation}
for stability in the norm $\lVert\,\cdot\,\rVert_h$.

If $f$ and $g$ are not zero, write $U^{n+1}=(I+\tau\Delta_h)U^n+\Delta t F^n$.
Under~\eqref{eq:CFL} (resp.~\eqref{eq:CFL-nu}),
\begin{equation}
\label{eq:l2stab-inhom}
\lVert U^n\rVert
\le
\lVert U^0\rVert
+\Delta t\sum_{j=0}^{n-1}\lVert F^j\rVert,
\end{equation}
where $F^j$ collects the source and boundary data at time $t_j$, and the
norm is $\ell^2(\mathcal{R}_h)$ in the uniform case and $\lVert\,\cdot\,\rVert_h$
in the nonuniform case.
\end{theorem}

\begin{proof}
On a uniform mesh the symbol of $\Delta_h$ is
$-(4/h^2)\sum_{\ell=1}^3\sin^2(\xi_\ell h/2)$ (cosine modes in the Neumann
case). Hence the spectrum lies in $[-12/h^2,0]$, with $\Lambda_h=12/h^2$.
Since $\Delta_h$ is symmetric, $\lVert I+\tau\Delta_h\rVert_{\ell^2\to\ell^2}$
equals the spectral radius of $I+\tau\Delta_h$. For $\lambda\in[-\Lambda_h,0]$,
the inequality $\lvert 1+\tau\lambda\rvert\le 1$ holds for all such $\lambda$
if and only if $0\le\tau\Lambda_h\le 2$, that is $\tau\le h^2/6$. This is
the equivalence of the three statements, and~\eqref{eq:l2stab} follows.

The same conclusion is reached by an energy identity, without Fourier
analysis. Because $-\Delta_h$ is symmetric positive semi-definite,
$\lVert\Delta_h V\rVert_{\ell^2}^2\le\Lambda_h\bigl(-(\Delta_h V,V)_{\ell^2}\bigr)$
for every $V$. Expanding~\eqref{eq:euler-hom},
\begin{equation}
\begin{aligned}
\lVert U^{n+1}\rVert_{\ell^2}^2
&=\lVert U^n\rVert_{\ell^2}^2
+2\tau(\Delta_h U^n,U^n)_{\ell^2}
+\tau^2\lVert\Delta_h U^n\rVert_{\ell^2}^2 \\
&\le
\lVert U^n\rVert_{\ell^2}^2
+2\tau(\Delta_h U^n,U^n)_{\ell^2}
+\tau^2\Lambda_h\bigl(-(\Delta_h U^n,U^n)_{\ell^2}\bigr) \\
&=\lVert U^n\rVert_{\ell^2}^2
+\bigl(2\tau-\tau^2\Lambda_h\bigr)(\Delta_h U^n,U^n)_{\ell^2}.
\end{aligned}
\end{equation}
The inner product $(\Delta_h U^n,U^n)_{\ell^2}$ is nonpositive. Hence
$\lVert U^{n+1}\rVert_{\ell^2}\le\lVert U^n\rVert_{\ell^2}$ as soon as
$2\tau-\tau^2\Lambda_h\ge 0$, that is $\tau\le 2/\Lambda_h=h^2/6$.

On a nonuniform mesh the operator defined by~\eqref{eq:dxx} is not symmetric
in the unweighted $\ell^2$ product, but it is symmetric in~\eqref{eq:weighted}.
The diagonal coefficient of $\delta_{xx}$ is $-2/(h_i h_{i-1})$, so the
Gershgorin radius of $-\delta_{xx}$ is at most $4/h_{\min}^2$. Summing the
three directions gives $\Lambda_h\le 12/h_{\min}^2$. Repeating the energy
estimate in the inner product $(\,\cdot\,,\,\cdot\,)_h$ yields~\eqref{eq:CFL-nu}.

In the inhomogeneous case, $\lVert I+\tau\Delta_h\rVert\le 1$ under the CFL
condition, therefore $\lVert U^{n+1}\rVert\le\lVert U^n\rVert+\Delta t\lVert F^n\rVert$.
Summing the resulting telescoping inequality yields~\eqref{eq:l2stab-inhom}.
\end{proof}

\begin{remark}
The threshold $h^2/6$ is the three-dimensional analogue of the classical
restrictions $k\Delta t\le h^2/2$ in one dimension and
$k\Delta t\le h^2/4$ in two dimensions, corresponding to a stencil with
$2d+1$ points. It is a restriction on forward Euler applied to the discrete
Laplacian; it is not a hyperbolic CFL condition of the form
$\Delta t\le C h$. Stability of a consistent scheme does not by itself imply
convergence of any particular order: an error estimate requires a truncation
analysis, which is not claimed here.

On the unit ball the same energy argument applies to the discrete spherical
Laplacian in the weighted $\ell^2$ product associated with the volume element
$\rho^2\sin\theta\,\dd\rho\,\dd\theta\,\dd\phi$. The constant $6$ is then
replaced by a constant that depends on the metric coefficients, and
$h_{\min}$ is the smallest radial or angular collocation spacing (the
restriction is most severe near the origin).
\end{remark}

%==============================================================================
\section{Numerical experiments}
\label{sec:numerics}

This section specifies the tests used to assess the scheme of
Section~\ref{sec:theory} on two standard three-dimensional geometries: the
unit cube and the unit ball. The cube study in Table~\ref{tab:cube} is the
explicit update~\eqref{eq:euler} with homogeneous Neumann closure, run with
$\Delta t$ at the CFL bound of Theorem~\ref{thm:scheme}; reported CPU times
are wall-clock times for the time loop. Linear finite-element reference
solutions, when used, are computed with piecewise affine tetrahedral elements
and the same time step as the \dziti{} scheme, so that the comparison
isolates the spatial discretisation. Errors are measured at the final time
$T$ in the discrete $\ell^2$ and $\ell^\infty$ norms on the collocation grid
$\mathcal{R}_h$,
\begin{equation}
\label{eq:errors}
e_2(h):=\Biggl(h^3\sum_{x_\alpha\in\mathcal{R}_h}
\bigl\lvert u(T,x_\alpha)-u_h(T,x_\alpha)\bigr\rvert^2\Biggr)^{1/2},
\qquad
e_\infty(h):=\max_{x_\alpha\in\mathcal{R}_h}
\bigl\lvert u(T,x_\alpha)-u_h(T,x_\alpha)\bigr\rvert.
\end{equation}
The observed order is
$p:=\log\bigl(e(h)/e(h/2)\bigr)/\log 2$. CPU times are wall-clock times for
the full time loop, excluding the (once-and-for-all) assembly of
$(\lambda_i,\psi_i(r_i))$.

\subsection{Detection of a breaking time: Burgers' equation}
\label{sec:burgers}

Before the three-dimensional tests we record a one-dimensional illustration
of the method's sensitivity to loss of regularity, which motivated its
original development \cite{Bsiss2016,Malek2022}. Consider the inviscid
Burgers equation
\begin{equation}
\label{eq:burgers}
\partial_t u+\partial_x\Bigl(\frac{u^2}{2}\Bigr)=0,
\qquad t>0,\quad x\in\mathbb{R},
\qquad
u(0,x)=\exp(-x^2/2).
\end{equation}
A classical solution exists only up to the breaking time
\begin{equation}
\label{eq:tstar}
t^\star
=-\frac{1}{\displaystyle\inf_{x\in\mathbb{R}}u_0'(x)}.
\end{equation}
Here $u_0'(x)=-x\exp(-x^2/2)$, so the infimum equals $u_0'(1)=-e^{-1/2}$ and
$t^\star=\sqrt{e}\simeq 1.648721$. Beyond $t^\star$ the entropy solution
develops a shock and the classical solution ceases to exist
\cite[Chapter~11]{LeVeque2002}. Figure~\ref{fig:burgers} shows the \dziti{}
approximation at several times, including a discrete time close to
$t^\star$. With $\Delta t=0.00915$ ($180$ time steps) the computed profile
ceases to remain single-valued at the theoretically predicted instant, after
a runtime of about six seconds. The same breaking time is obtained by the
Lax--Friedrichs scheme in \cite{Yoon2008}, at a higher mesh cost. This test
is not a substitute for the parabolic analysis of Section~\ref{sec:theory};
it illustrates why a basis built from nascent Dirac masses is of interest
when the solution may leave $H^1$.

\begin{figure}[htbp]
\centering
\includegraphics[width=0.85\textwidth,height=1.8in]{Burger}
\caption{Inviscid Burgers equation~\eqref{eq:burgers}: \dziti{} approximation
at $t=0.0915$, $t=1.0986$ and $t=1.6485\simeq t^\star=\sqrt{e}$.}
\label{fig:burgers}
\end{figure}

\subsection{Transient diffusion in the unit cube}
\label{sec:cube}

Let $\Omega=(0,1)^3$ and $k\equiv 1$. The manufactured solution
\begin{equation}
\label{eq:uex-cube}
u_{\mathrm{ex}}(x,y,z,t)
=\cos(\pi x)\cos(\pi y)\cos(\pi z)\,e^{-3\pi^2 t}
\end{equation}
satisfies~\eqref{eq:heat} with $f=0$, homogeneous Neumann boundary conditions
$g=0$, and initial datum $u_0(x,y,z)=\cos(\pi x)\cos(\pi y)\cos(\pi z)$. The
scheme~\eqref{eq:euler} is run up to $T=0.1$ on a sequence of collocation
grids with $m+1\in\{8,16,32,64\}$ nodes per direction, with $\Delta t$
obeying~\eqref{eq:CFL}, taking $\Delta t=T/N_t$ with $N_t$ the smallest
integer such that $k\Delta t\le h^2/6$. Table~\ref{tab:cube} reports $e_2$,
$e_\infty$, the observed order $p_2$ and the CPU time. The observed $L^2$
orders approach $2$, in agreement with second-order consistency of the
seven-point Laplacian on a uniform mesh.

A second cube test uses a non-separable source,
\begin{equation}
\label{eq:uex-cube2}
u_{\mathrm{ex}}(x,y,z,t)
=\exp\bigl(-(x^2+y^2+z^2)\bigr)\,e^{-t},
\end{equation}
from which $f$ and the Neumann datum $g=k\partial_n u_{\mathrm{ex}}$ are
computed analytically. This checks the inhomogeneous boundary closure.

\begin{table}[htbp]
\centering
\caption{Unit cube, manufactured solution~\eqref{eq:uex-cube} at $T=0.1$:
\dziti{} errors, observed $L^2$ order $p_2$ and CPU time.}
\label{tab:cube}
\begin{tabular}{crrrrr}
\toprule
$m+1$ & $N_h=(m+1)^3$ & $e_2$ & $e_\infty$ & $p_2$ & CPU (s) \\
\midrule
$8$  & $512$     & $2.58\times 10^{-3}$ & $5.01\times 10^{-3}$ & ---   & $0.001$ \\
$16$ & $4096$    & $4.78\times 10^{-4}$ & $1.12\times 10^{-3}$ & $2.43$ & $0.008$ \\
$32$ & $32768$   & $1.02\times 10^{-4}$ & $2.62\times 10^{-4}$ & $2.23$ & $0.32$ \\
$64$ & $262144$  & $2.35\times 10^{-5}$ & $6.35\times 10^{-5}$ & $2.11$ & $8.73$ \\
\bottomrule
\end{tabular}
\end{table}

\subsection{Transient diffusion in the unit ball}
\label{sec:ball}

Let $\Omega=B(0,1)\subset\mathbb{R}^3$ and $k\equiv 1$. In spherical
coordinates we take the radial manufactured solution
\begin{equation}
\label{eq:uex-ball}
u_{\mathrm{ex}}(\rho,t)
=\bigl(1-\tfrac13\rho^2\bigr)e^{-t},
\end{equation}
for which
\begin{equation}
\label{eq:fg-ball}
f(\rho,t)
=\bigl(-1+\tfrac13\rho^2+2\bigr)e^{-t}
=\bigl(1+\tfrac13\rho^2\bigr)e^{-t},
\qquad
g(\theta,\phi,t)
=k\partial_\rho u_{\mathrm{ex}}(1,t)
=-\tfrac23 e^{-t}.
\end{equation}
(The identity $\Delta(\rho^2)=6$ was used.) The spherical-coordinate
\dziti{} scheme is compared with the chord-sweeping variant and with linear
finite elements on a tetrahedral mesh of the ball. Errors are computed as
in~\eqref{eq:errors}, with the spherical volume weights included in $e_2$.
A second, non-radial test is obtained by transporting~\eqref{eq:uex-cube2}
to the ball, which breaks spherical symmetry and exercises the angular
differences in~\eqref{eq:sph-lap}.

In both geometries we additionally record the discrete energy
$E^n:=\lVert U^n\rVert_{\ell^2(\mathcal{R}_h)}^2$ for the homogeneous
Neumann problem, in order to check the decay predicted by
Theorem~\ref{thm:semi} and the CFL restriction of Theorem~\ref{thm:scheme}.

\begin{table}[htbp]
\centering
\caption{Unit ball, manufactured solution~\eqref{eq:uex-ball} at $T=0.1$.
Spherical-coordinate \dziti{} scheme versus chord sweeping and linear FEM.}
\label{tab:ball}
\begin{tabular}{llrrrr}
\toprule
Method & $N_h$ & $e_2$ & $e_\infty$ & $p_2$ & CPU (s) \\
\midrule
\dziti{} (spherical) &  &  &  &  &  \\
\dziti{} (chords)    &  &  &  &  &  \\
linear FEM           &  &  &  &  &  \\
\bottomrule
\end{tabular}
\end{table}

%==============================================================================
\section{Conclusion}
\label{sec:concl}

We have given a mathematically consistent account of the \dziti{} method for
three-dimensional transient diffusion with Neumann boundary conditions. The
one-dimensional generators are compactly supported mollifiers; their Gram
matrix is tridiagonal, which reduces orthonormalisation to a two-term
recurrence and produces an orthonormal basis of a finite-dimensional subspace
$V_h\subset L^2$. The roots of that basis define the reconstruction $f_h$
of Lemma~\ref{lem:orthonormal} and a quadrature that avoids numerical
integration rules of Gaussian type. By that lemma, $f_h$ is the best
least-squares approximation in $V_h$, and it is compared with the interpolant
$I_h f=\sum_{i=1}^{m+1}f(x_i)\,\varphi_i(x)$. Theorem~\ref{thm:mollify} gives
the second-order mollification bound on $\eta_h*f$. Theorem~\ref{thm:error}
shows that no comparable bound can hold for $\lVert f-f_h\rVert_{L^2}$, because
$V_h$ does not contain constants.

The semi-discrete Galerkin scheme is unconditionally energy-stable. The fully
discrete explicit Euler scheme on a uniform Cartesian mesh is stable in
$\ell^2$ if and only if $k\Delta t\le h^2/6$; on a nonuniform mesh the same
threshold with $h$ replaced by $h_{\min}$ is sufficient in the weighted
discrete $L^2$ norm. The scheme requires no mass-matrix inversion. The same
building blocks tensorise to the unit cube and extend to the unit ball by
spherical coordinates or by chord sweeping.

Several questions remain. The Neumann closure on the ball, the treatment of
variable or anisotropic conductivities, and a fully discrete error analysis
in $L^2(0,T;H^1(\Omega))$ are left to future work. On the computational side,
a comparison with PUFEM \cite{Malek2019,Mohamed2013} at equal accuracy,
including problems with steep boundary layers, would quantify the cost
reduction permitted by the explicit \dziti{} update.

\section*{Acknowledgements}

The author thanks the developers of the \dziti{} method, in particular
C.~Ziti and L.~Bsiss, whose original construction this analysis builds upon.

\bibliographystyle{elsarticle-num}
\bibliography{bibliography}

\end{document}
